% ============================================================
%  C: Como Programar -- 6ª Edição | Deitel & Deitel
%  Capítulo 12 -- Estruturas de Dados em C
%  EXERCÍCIOS (12.6 -- 12.25)
%  Abordagem incremental: Caps. 1 a 12
% ============================================================
\documentclass[12pt,a4paper]{article}
\usepackage[utf8]{inputenc}
\usepackage[T1]{fontenc}
\usepackage[brazil]{babel}
\usepackage{geometry}
\geometry{top=2.5cm,bottom=2.5cm,left=3cm,right=3cm}
\usepackage{parskip}\setlength{\parskip}{6pt}\setlength{\parindent}{0pt}
\usepackage{xcolor,tcolorbox,enumitem,listings,fancyhdr,titlesec,hyperref,booktabs,array,amsmath}
\tcbuselibrary{skins,breakable}

\definecolor{deitelblue}{RGB}{0,70,127}
\definecolor{deitelgray}{RGB}{240,242,245}
\definecolor{deitelgreen}{RGB}{0,110,60}
\definecolor{deitellightblue}{RGB}{220,235,250}
\definecolor{deitelred}{RGB}{160,20,20}
\definecolor{deitellorange}{RGB}{200,90,0}

\newtcolorbox{boaprática}[1][]{colback=deitellightblue,colframe=deitelblue,
  fonttitle=\bfseries\small,title={Boa prática de programação},breakable,#1}
\newtcolorbox{erroco}[1][]{colback=red!5,colframe=deitelred,
  fonttitle=\bfseries\small,title={Erro comum de programação},breakable,#1}
\newtcolorbox{respostabox}[1][]{colback=deitelgray,colframe=deitelblue!60,
  fonttitle=\bfseries\small\color{deitelblue},title={Resposta detalhada},breakable,#1}
\newtcolorbox{enunciadoBox}[1][]{colback=white,colframe=deitelblue,
  fonttitle=\bfseries,title={#1},breakable}
\newtcolorbox{saida}[1][]{colback=black!88,colframe=deitelblue,
  fontupper=\ttfamily\small\color{green!70},
  title={\small\color{white}Saída do programa},breakable,#1}

\setlist[enumerate,1]{label=\textbf{\alph*)},leftmargin=2em}
\setlist[itemize]{leftmargin=2em}

\lstset{
  basicstyle=\ttfamily\small,backgroundcolor=\color{deitelgray},
  frame=single,framerule=0.4pt,rulecolor=\color{deitelblue!40},
  breaklines=true,showstringspaces=false,
  keywordstyle=\color{deitelblue}\bfseries,
  commentstyle=\color{deitelgreen}\itshape,
  stringstyle=\color{deitelred},
  language=C,numbers=left,numberstyle=\tiny\color{gray},
  xleftmargin=18pt,xrightmargin=5pt,stepnumber=1,
}

\pagestyle{fancy}\fancyhf{}
\fancyhead[L]{\small\color{deitelblue}\textbf{C: Como Programar -- 6ª ed. | Deitel \& Deitel}}
\fancyhead[R]{\small\color{deitelblue}Cap. 12 -- Exercícios}
\fancyfoot[C]{\small\thepage}
\renewcommand{\headrulewidth}{0.4pt}

\titleformat{\section}[block]{\large\bfseries\color{deitelblue}}{\thesection.}{0.5em}{}[\titlerule]
\titleformat{\subsection}[block]{\normalsize\bfseries\color{deitelblue!80}}{}{}{}

\hypersetup{colorlinks=true,linkcolor=deitelblue}

\begin{document}

\begin{titlepage}
  \centering\vspace*{2cm}
  {\color{deitelblue}\rule{\linewidth}{2pt}}\\[0.4cm]
  {\huge\bfseries\color{deitelblue}C: Como Programar}\\[0.2cm]
  {\large\color{deitelblue}6ª Edição $\cdot$ Paul Deitel \& Harvey Deitel}\\[0.2cm]
  {\color{deitelblue}\rule{\linewidth}{2pt}}\\[1cm]
  {\LARGE\bfseries Capítulo 12}\\[0.3cm]
  {\Large\color{deitelblue!80}Estruturas de Dados em C}\\[0.8cm]
  \begin{tcolorbox}[colback=deitellightblue,colframe=deitelblue,width=0.85\linewidth]
    \centering{\large\bfseries Exercícios (12.6--12.25)}\\[0.2cm]
    {\normalsize Resolvidos passo a passo, com código completo}
  \end{tcolorbox}
\end{titlepage}

\tableofcontents\newpage

% ============================================================
\section{Exercício 12.6 -- Concatenação de Listas}
% ============================================================

\begin{respostabox}
\begin{lstlisting}
/* Ex12.6: concatenar_listas.c */
#include <stdio.h>
#include <stdlib.h>

typedef struct charNode {
    char data;
    struct charNode *nextPtr;
} CharNode;

void concatenate( CharNode **firstPtr, CharNode *secondPtr )
{
    if ( *firstPtr == NULL ) {
        *firstPtr = secondPtr;
        return;
    }
    /* Avanca ate o ultimo no da primeira lista */
    CharNode *p = *firstPtr;
    while ( p->nextPtr != NULL ) {
        p = p->nextPtr;
    }
    p->nextPtr = secondPtr;  /* engata a segunda lista */
}

void inserir( CharNode **head, char c )
{
    CharNode *novo = malloc( sizeof( CharNode ) );
    novo->data = c;
    novo->nextPtr = NULL;

    if ( *head == NULL ) {
        *head = novo;
    } else {
        CharNode *p = *head;
        while ( p->nextPtr != NULL ) p = p->nextPtr;
        p->nextPtr = novo;
    }
}

void imprimir( CharNode *head )
{
    while ( head != NULL ) {
        printf( "%c ", head->data );
        head = head->nextPtr;
    }
    printf( "\n" );
}

int main( void )
{
    CharNode *l1 = NULL, *l2 = NULL;
    inserir( &l1, 'A' );  inserir( &l1, 'B' );  inserir( &l1, 'C' );
    inserir( &l2, 'X' );  inserir( &l2, 'Y' );  inserir( &l2, 'Z' );

    printf( "L1: " );  imprimir( l1 );
    printf( "L2: " );  imprimir( l2 );

    concatenate( &l1, l2 );

    printf( "Concatenada: " );  imprimir( l1 );

    return 0;
}
\end{lstlisting}

\textbf{Conceito-chave:} concatenação em $O(n)$ -- precisamos percorrer a primeira
lista até o fim para então engatar a segunda. Listas duplamente encadeadas com
ponteiro para o último nó tornariam isso $O(1)$.
\end{respostabox}

\newpage

% ============================================================
\section{Exercício 12.7 -- Mescla de Listas Ordenadas}
% ============================================================

\begin{respostabox}
\begin{lstlisting}
/* Ex12.7: mesclar_ordenadas.c */
#include <stdio.h>
#include <stdlib.h>

typedef struct intNode {
    int data;
    struct intNode *nextPtr;
} IntNode;

IntNode* merge( IntNode *a, IntNode *b )
{
    IntNode dummy;        /* no sentinela para simplificar */
    IntNode *tail = &dummy;
    dummy.nextPtr = NULL;

    while ( a != NULL && b != NULL ) {
        if ( a->data <= b->data ) {
            tail->nextPtr = a;
            a = a->nextPtr;
        } else {
            tail->nextPtr = b;
            b = b->nextPtr;
        }
        tail = tail->nextPtr;
    }
    /* Conecta o que sobrou */
    tail->nextPtr = ( a != NULL ) ? a : b;

    return dummy.nextPtr;
}
\end{lstlisting}

\textbf{Algoritmo:} compare as cabeças, escolha a menor, avance. Quando uma lista
acabar, conecte a outra. $O(n+m)$.

Este é o ``merge'' do \textit{merge sort} -- uma das ideias mais elegantes da
ciência da computação.
\end{respostabox}

% ============================================================
\section{Exercício 12.8 -- Inserção Ordenada de 25 Inteiros}
% ============================================================

\begin{respostabox}
\begin{lstlisting}
/* Ex12.8: insercao_ordenada.c */
#include <stdio.h>
#include <stdlib.h>
#include <time.h>

typedef struct node {
    int data;
    struct node *nextPtr;
} Node;

void inserirOrdenado( Node **head, int valor )
{
    Node *novo = malloc( sizeof( Node ) );
    novo->data = valor;

    Node *prev = NULL, *cur = *head;
    while ( cur != NULL && cur->data < valor ) {
        prev = cur;  cur = cur->nextPtr;
    }
    novo->nextPtr = cur;
    if ( prev == NULL ) *head = novo;
    else prev->nextPtr = novo;
}

int main( void )
{
    Node *lista = NULL;
    int i, valor, soma = 0;

    srand( time( NULL ) );

    for ( i = 0; i < 25; ++i ) {
        valor = rand() % 101;  /* 0 a 100 */
        inserirOrdenado( &lista, valor );
        soma += valor;
    }

    printf( "Lista ordenada:\n" );
    Node *p = lista;
    while ( p != NULL ) { printf( "%d ", p->data ); p = p->nextPtr; }
    printf( "\nSoma:   %d\nMedia: %.2f\n", soma, soma / 25.0 );

    return 0;
}
\end{lstlisting}
\end{respostabox}

\newpage

% ============================================================
\section{Exercício 12.9 -- Lista Invertida}
% ============================================================

\begin{respostabox}
\begin{lstlisting}
/* Ex12.9: inverter_lista.c
   Cria lista de 10 chars e gera copia invertida */
#include <stdio.h>
#include <stdlib.h>

typedef struct cnode {
    char data;
    struct cnode *nextPtr;
} CNode;

void inserirFim( CNode **head, char c )
{
    CNode *novo = malloc( sizeof( CNode ) );
    novo->data = c;  novo->nextPtr = NULL;
    if ( *head == NULL ) { *head = novo; return; }
    CNode *p = *head;
    while ( p->nextPtr != NULL ) p = p->nextPtr;
    p->nextPtr = novo;
}

/* Insere SEMPRE no inicio -> inverte naturalmente */
void inserirInicio( CNode **head, char c )
{
    CNode *novo = malloc( sizeof( CNode ) );
    novo->data = c;
    novo->nextPtr = *head;
    *head = novo;
}

int main( void )
{
    CNode *original = NULL, *inversa = NULL;
    char chars[] = "ABCDEFGHIJ";
    int i;

    for ( i = 0; i < 10; ++i ) inserirFim( &original, chars[ i ] );

    /* Percorre original e insere no inicio da inversa */
    CNode *p = original;
    while ( p != NULL ) {
        inserirInicio( &inversa, p->data );
        p = p->nextPtr;
    }

    printf( "Original: " );
    p = original; while ( p != NULL ) { putchar( p->data ); p = p->nextPtr; }
    printf( "\nInversa:  " );
    p = inversa;  while ( p != NULL ) { putchar( p->data ); p = p->nextPtr; }
    printf( "\n" );
    return 0;
}
\end{lstlisting}

\textbf{Truque elegante:} ao percorrer a original do início ao fim e \textit{inserir
no início} da nova lista, obtemos a ordem invertida automaticamente.
\end{respostabox}

\newpage

% ============================================================
\section{Exercício 12.10 -- Inverter Linha de Texto}
% ============================================================

\begin{respostabox}
\begin{lstlisting}
/* Ex12.10: inverter_texto.c
   Usa pilha para inverter caracteres de uma linha */
#include <stdio.h>
#include <stdlib.h>

typedef struct snode {
    char data;
    struct snode *nextPtr;
} SNode;

void push( SNode **top, char c )
{
    SNode *novo = malloc( sizeof( SNode ) );
    novo->data = c;
    novo->nextPtr = *top;
    *top = novo;
}

char pop( SNode **top )
{
    SNode *temp = *top;
    char c = temp->data;
    *top = temp->nextPtr;
    free( temp );
    return c;
}

int main( void )
{
    SNode *pilha = NULL;
    int c;

    printf( "Digite uma linha: " );
    while ( ( c = getchar() ) != '\n' && c != EOF ) {
        push( &pilha, c );
    }

    printf( "Invertida:  " );
    while ( pilha != NULL ) {
        putchar( pop( &pilha ) );
    }
    putchar( '\n' );
    return 0;
}
\end{lstlisting}

\textbf{Por que pilha funciona aqui?} Empilhamos na ordem em que lemos e
desempilhamos na ordem inversa -- exatamente o que precisamos. LIFO em ação.
\end{respostabox}

% ============================================================
\section{Exercício 12.11 -- Testador de Palíndromos}
% ============================================================

\begin{respostabox}
\textbf{Estratégia:} usar uma pilha e uma fila simultaneamente. Cada caractere
significativo (não-espaço, não-pontuação) é inserido em ambos. No fim,
desempilhar e desenfileirar comparando -- se diferentes em algum ponto, não é
palíndromo.

\begin{lstlisting}
/* Ex12.11: palindromo.c (esqueleto) */
#include <stdio.h>
#include <ctype.h>

/* (defina pilha e fila para char) */

int main( void )
{
    SNode *pilha = NULL;
    QNode *frenteFila = NULL, *finalFila = NULL;
    int c;

    printf( "Digite uma frase: " );
    while ( ( c = getchar() ) != '\n' && c != EOF ) {
        if ( isalpha( c ) ) {
            char ch = tolower( c );
            push( &pilha, ch );
            enqueue( &frenteFila, &finalFila, ch );
        }
    }

    int palindromo = 1;
    while ( pilha != NULL && frenteFila != NULL ) {
        if ( pop( &pilha ) != dequeue( &frenteFila ) ) {
            palindromo = 0;
            break;
        }
    }

    printf( palindromo ? "Eh palindromo!\n"
                       : "Nao eh palindromo.\n" );
    return 0;
}
\end{lstlisting}

\textbf{Exemplos:} ``A man a plan a canal Panama'', ``A base seba'', ``Anotaram a
data da maratona''. Espaços e pontuação são ignorados.
\end{respostabox}

\newpage

% ============================================================
\section{Exercício 12.12 -- Conversão Infixo $\to$ Pós-fixo}
% ============================================================

\begin{respostabox}
\textbf{Algoritmo Shunting Yard (Dijkstra), versão simplificada para dígitos
únicos:}

\begin{lstlisting}
/* Ex12.12: infix_to_postfix.c -- estrutura central */
#include <stdio.h>
#include <stdlib.h>
#include <string.h>

typedef struct stackNode {
    char data;
    struct stackNode *nextPtr;
} StackNode;
typedef StackNode *StackNodePtr;

void push( StackNodePtr *topPtr, char value );
char pop( StackNodePtr *topPtr );
char stackTop( StackNodePtr topPtr );
int  isEmpty( StackNodePtr topPtr );
int  isOperator( char c );
int  precedence( char op1, char op2 );

void convertToPostfix( char infix[], char postfix[] )
{
    StackNodePtr stack = NULL;
    int infixIndex = 0, postfixIndex = 0;
    char c;

    push( &stack, '(' );           /* etapa 1 */
    strcat( infix, ")" );          /* etapa 2 */

    while ( !isEmpty( stack ) ) {
        c = infix[ infixIndex ];

        if ( c >= '0' && c <= '9' ) {
            postfix[ postfixIndex++ ] = c;
        }
        else if ( c == '(' ) {
            push( &stack, c );
        }
        else if ( isOperator( c ) ) {
            while ( !isEmpty( stack ) &&
                    isOperator( stackTop( stack ) ) &&
                    precedence( stackTop( stack ), c ) >= 0 ) {
                postfix[ postfixIndex++ ] = pop( &stack );
            }
            push( &stack, c );
        }
        else if ( c == ')' ) {
            while ( stackTop( stack ) != '(' ) {
                postfix[ postfixIndex++ ] = pop( &stack );
            }
            pop( &stack );  /* descarta '(' */
        }
        ++infixIndex;
    }
    postfix[ postfixIndex ] = '\0';
}

int isOperator( char c )
{
    return c == '+' || c == '-' || c == '*' ||
           c == '/' || c == '^' || c == '%';
}

int precedence( char op1, char op2 )
{
    /* Retorna 1 se prec(op1) > prec(op2)
              0 se iguais
             -1 se prec(op1) < prec(op2)  */
    int p1, p2;
    if ( op1 == '^' )                          p1 = 3;
    else if ( op1 == '*' || op1 == '/' ||
              op1 == '%' )                     p1 = 2;
    else                                       p1 = 1;
    if ( op2 == '^' )                          p2 = 3;
    else if ( op2 == '*' || op2 == '/' ||
              op2 == '%' )                     p2 = 2;
    else                                       p2 = 1;
    if ( p1 == p2 ) return 0;
    return p1 > p2 ? 1 : -1;
}
\end{lstlisting}

\textbf{Exemplo:} \texttt{(6 + 2) * 5 - 8 / 4} $\to$ \texttt{6 2 + 5 * 8 4 / -}

\textbf{Por que pilha?} Operadores precisam ``esperar'' seus operandos. A pilha
guarda operadores até que a expressão esteja pronta para emiti-los.
\end{respostabox}

\newpage

% ============================================================
\section{Exercício 12.13 -- Avaliador Pós-fixo}
% ============================================================

\begin{respostabox}
\begin{lstlisting}
/* Ex12.13: postfix_eval.c */
#include <stdio.h>
#include <stdlib.h>
#include <string.h>
#include <math.h>

typedef struct stackNode {
    int data;
    struct stackNode *nextPtr;
} StackNode;
typedef StackNode *StackNodePtr;

void push( StackNodePtr *topPtr, int value );
int  pop( StackNodePtr *topPtr );
int  isEmpty( StackNodePtr topPtr );
int  calculate( int op1, int op2, char op );

int evaluatePostfixExpression( char *expr )
{
    StackNodePtr stack = NULL;
    int i = 0;
    char c;

    strcat( expr, "\0" );  /* etapa 1 */

    while ( ( c = expr[ i ] ) != '\0' ) {
        if ( c >= '0' && c <= '9' ) {
            push( &stack, c - '0' );  /* converte char -> int */
        }
        else if ( c == '+' || c == '-' || c == '*' ||
                  c == '/' || c == '^' || c == '%' ) {
            int x = pop( &stack );
            int y = pop( &stack );
            push( &stack, calculate( y, x, c ) );
        }
        ++i;
    }

    return pop( &stack );
}

int calculate( int op1, int op2, char op )
{
    switch ( op ) {
        case '+': return op1 + op2;
        case '-': return op1 - op2;
        case '*': return op1 * op2;
        case '/': return op1 / op2;
        case '%': return op1 % op2;
        case '^': return ( int )pow( op1, op2 );
    }
    return 0;
}
\end{lstlisting}

\textbf{Trace de \texttt{6 2 + 5 * 8 4 / -}:}
\begin{enumerate}[noitemsep]
  \item \texttt{6} $\to$ pilha [6]
  \item \texttt{2} $\to$ pilha [6,2]
  \item \texttt{+}: pop 2, pop 6; push 6+2 $\to$ [8]
  \item \texttt{5} $\to$ [8,5]
  \item \texttt{*}: pop 5, pop 8; push 8*5 $\to$ [40]
  \item \texttt{8} $\to$ [40,8]
  \item \texttt{4} $\to$ [40,8,4]
  \item \texttt{/}: pop 4, pop 8; push 8/4 $\to$ [40,2]
  \item \texttt{-}: pop 2, pop 40; push 40-2 $\to$ [38]
\end{enumerate}
Resultado: \textbf{38}.
\end{respostabox}

% ============================================================
\section{Exercício 12.14 -- Operandos com Múltiplos Dígitos}
% ============================================================

\begin{respostabox}
\textbf{Modificação:} agora os operandos podem ter vários dígitos. Como
distinguir o fim de um número? Usar um \textit{delimitador}, normalmente espaço.

\begin{lstlisting}
/* Loop principal modificado */
i = 0;
while ( expr[ i ] != '\0' ) {
    char c = expr[ i ];

    if ( c >= '0' && c <= '9' ) {
        int valor = 0;
        /* Acumula digitos enquanto for digito */
        while ( expr[ i ] >= '0' && expr[ i ] <= '9' ) {
            valor = valor * 10 + ( expr[ i ] - '0' );
            ++i;
        }
        push( &stack, valor );
        continue;  /* nao incrementa i de novo */
    }
    /* ... resto igual ... */
    ++i;
}
\end{lstlisting}

\textbf{Princípio:} construir o número dígito a dígito multiplicando o acumulador
por 10. Esta é a base de todos os parsers numéricos -- de \texttt{atoi} aos
\textit{lexers} de compiladores.
\end{respostabox}

\newpage

% ============================================================
\section{Exercício 12.15 -- Simulação de Supermercado}
% ============================================================

\begin{respostabox}
\begin{lstlisting}
/* Ex12.15: supermercado.c -- esqueleto */
#include <stdio.h>
#include <stdlib.h>
#include <time.h>

typedef struct cliente {
    int chegada;          /* minuto em que chegou na fila */
    struct cliente *next;
} Cliente;

Cliente *cabeca = NULL, *cauda = NULL;

void enqueue( int t );
int  dequeue( void );
int  tamanhoFila( void );

int main( void )
{
    srand( time( NULL ) );

    int prox_chegada = 1 + rand() % 4;
    int caixa_livre_em = 0;
    int max_fila = 0, max_espera = 0;

    int t;
    for ( t = 1; t <= 720; ++t ) {
        /* Chegada */
        if ( t == prox_chegada ) {
            enqueue( t );
            int sz = tamanhoFila();
            if ( sz > max_fila ) max_fila = sz;
            prox_chegada = t + 1 + rand() % 4;
        }
        /* Caixa livre? */
        if ( t >= caixa_livre_em && cabeca != NULL ) {
            int chegou = dequeue();
            int espera = t - chegou;
            if ( espera > max_espera ) max_espera = espera;
            caixa_livre_em = t + 1 + rand() % 4;
        }
    }

    printf( "Maxima fila:  %d\n", max_fila );
    printf( "Espera maxima: %d minutos\n", max_espera );
    return 0;
}
\end{lstlisting}

\textbf{Respostas típicas (variam por aleatoriedade):}
\begin{itemize}[noitemsep]
  \item \textbf{(a) Fila máxima:} tipicamente 2--5 clientes.
  \item \textbf{(b) Espera máxima:} tipicamente 5--15 minutos.
  \item \textbf{(c) Se chegada mudar para 1--3 min:} taxa de chegada ($\sim$2 min/cliente)
  passa a ser menor que a de atendimento ($\sim$2.5 min/cliente), causando
  \textbf{crescimento sem limites} da fila ao longo do dia.
\end{itemize}

\textbf{Lição:} este é um exemplo de \textit{teoria das filas} -- ramo da matemática
aplicada à engenharia de tráfego, redes, call centers, etc.
\end{respostabox}

% ============================================================
\section{Exercício 12.16 -- Permitir Duplicatas}
% ============================================================

\begin{respostabox}
\textbf{Modificação simples:} na função \texttt{insertNode}, alterar a condição
de comparação para aceitar iguais em um dos lados:

\begin{lstlisting}
/* Original: descarta duplicatas
   if ( valor < no->data ) ... esquerda
   else if ( valor > no->data ) ... direita
   else ; // duplicata -- ignora

   Modificado: aceita duplicatas, vai para a direita: */
if ( valor < no->data ) {
    /* insere a esquerda */
} else {  /* valor >= no->data */
    /* insere a direita */
}
\end{lstlisting}

\textbf{Consequência:} a travessia in-order ainda produz ordem crescente, mas com
duplicatas adjacentes. Outras operações (busca, exclusão) funcionam normalmente.
\end{respostabox}

\newpage

% ============================================================
\section{Exercício 12.17 -- Árvore de Strings}
% ============================================================

\begin{respostabox}
\begin{lstlisting}
/* Ex12.17: arvore_strings.c -- esqueleto */
#include <stdio.h>
#include <stdlib.h>
#include <string.h>

typedef struct treeNode {
    char *str;
    struct treeNode *left, *right;
} TreeNode;

void insert( TreeNode **root, char *palavra )
{
    if ( *root == NULL ) {
        TreeNode *novo = malloc( sizeof( TreeNode ) );
        novo->str = malloc( strlen( palavra ) + 1 );
        strcpy( novo->str, palavra );
        novo->left = novo->right = NULL;
        *root = novo;
        return;
    }
    int cmp = strcmp( palavra, (*root)->str );
    if      ( cmp < 0 ) insert( &(*root)->left,  palavra );
    else if ( cmp > 0 ) insert( &(*root)->right, palavra );
    /* iguais: ignora (sem duplicatas) */
}

int main( void )
{
    TreeNode *raiz = NULL;
    char texto[ 1000 ];
    char *tok;

    printf( "Digite uma linha de texto:\n" );
    fgets( texto, 1000, stdin );

    tok = strtok( texto, " \t\n.,;:!?" );
    while ( tok != NULL ) {
        insert( &raiz, tok );
        tok = strtok( NULL, " \t\n.,;:!?" );
    }

    /* Chamar inOrder, preOrder, postOrder */
    return 0;
}
\end{lstlisting}

\textbf{Nota:} a travessia in-order da árvore produz as palavras em ordem
alfabética -- excelente para criar índices de livros ou dicionários.
\end{respostabox}

% ============================================================
\section{Exercício 12.18 -- Eliminação com Array vs Árvore}
% ============================================================

\begin{respostabox}
\textbf{Eliminação com array:}
\begin{enumerate}[noitemsep]
  \item Para cada novo valor, percorrer todo o array procurando duplicata.
  \item Se não encontrar, inserir no final.
  \item Complexidade: $O(n)$ por inserção; total $O(n^2)$ para $n$ valores.
\end{enumerate}

\textbf{Eliminação com árvore binária de busca:}
\begin{enumerate}[noitemsep]
  \item Para cada novo valor, descer pela árvore.
  \item Se chegar a um nó vazio: inserir.
  \item Se encontrar igual: ignorar.
  \item Complexidade média: $O(\log n)$ por inserção; total $O(n \log n)$.
  \item Pior caso (árvore desbalanceada, dados já ordenados): $O(n^2)$.
\end{enumerate}

\textbf{Comparação:}
\begin{center}
\renewcommand{\arraystretch}{1.25}
\begin{tabular}{l c c}
\toprule
\textbf{Estrutura} & \textbf{Médio} & \textbf{Pior caso} \\ \midrule
Array            & $O(n^2)$        & $O(n^2)$ \\
Árvore BST       & $O(n \log n)$   & $O(n^2)$ \\
Árvore balanceada (AVL, RB) & $O(n \log n)$ & $O(n \log n)$ \\
Tabela hash      & $O(n)$          & $O(n^2)$ raro \\
\bottomrule
\end{tabular}
\end{center}

Para grandes conjuntos de dados, árvores balanceadas ou tabelas hash são as
escolhas profissionais.
\end{respostabox}

\newpage

% ============================================================
\section{Exercício 12.19 -- Profundidade da Árvore}
% ============================================================

\begin{respostabox}
\begin{lstlisting}
/* Ex12.19: profundidade.c */
int depth( const TreeNode *root )
{
    if ( root == NULL ) return 0;

    int d_esq = depth( root->left );
    int d_dir = depth( root->right );

    return 1 + ( d_esq > d_dir ? d_esq : d_dir );
}
\end{lstlisting}

\textbf{Análise:}
\begin{itemize}[noitemsep]
  \item Caso base: árvore vazia tem profundidade 0.
  \item Recursão: profundidade = 1 (este nó) + máximo das profundidades das subárvores.
\end{itemize}

\textbf{Para a árvore do Ex.~12.5:} a profundidade é 4 níveis (raiz 49, depois
28/83, depois 18/40/71/97, e por fim 11/19/32/44/69/72/92/99).
\end{respostabox}

% ============================================================
\section{Exercício 12.20 -- Impressão Recursiva Reversa}
% ============================================================

\begin{respostabox}
\begin{lstlisting}
/* Ex12.20: print_reverse.c */
void printListBackwards( const Node *p )
{
    if ( p == NULL ) return;

    printListBackwards( p->nextPtr );  /* chama primeiro */
    printf( "%d ", p->data );           /* imprime depois */
}
\end{lstlisting}

\textbf{Por que funciona?} Cada chamada recursiva mergulha mais fundo na lista
\textit{antes} de imprimir. Quando chega ao fim, retorna -- e a impressão
acontece na ordem reversa do empilhamento de chamadas. A pilha de chamadas
faz o trabalho da pilha explícita do Ex.~12.10.

\textbf{Cuidado:} para listas muito longas, isso pode estourar a pilha de
chamadas (\textit{stack overflow}). Para essas, prefira a versão iterativa
com pilha explícita.
\end{respostabox}

% ============================================================
\section{Exercício 12.21 -- Busca Recursiva}
% ============================================================

\begin{respostabox}
\begin{lstlisting}
/* Ex12.21: search_recursive.c */
Node* searchList( Node *head, int valor )
{
    if ( head == NULL )           return NULL;
    if ( head->data == valor )    return head;
    return searchList( head->nextPtr, valor );
}
\end{lstlisting}

\textbf{Três casos:}
\begin{enumerate}[noitemsep]
  \item Lista vazia: retorna \texttt{NULL} (não achou).
  \item Achou: retorna o ponteiro para o nó.
  \item Senão: busca recursivamente no resto.
\end{enumerate}

\textbf{Padrão:} esta é uma busca em lista em $O(n)$ -- mesma complexidade da
versão iterativa, mas usa pilha de chamadas. Tail-recursive: bons compiladores
otimizam para um loop interno.
\end{respostabox}

\newpage

% ============================================================
\section{Exercício 12.22 -- Exclusão em Árvore Binária}
% ============================================================

\begin{respostabox}
\textbf{Três casos da exclusão:}

\textbf{Caso 1 -- Nó folha:}
Simplesmente remove o nó e ajusta o pai para \texttt{NULL}.

\textbf{Caso 2 -- Nó com um filho:}
Substitui o nó pelo seu único filho.

\textbf{Caso 3 -- Nó com dois filhos:}
Procura o \textit{substituto} (predecessor in-order: o maior valor da subárvore
esquerda). Copia o valor do substituto para o nó atual e remove o substituto
(que é caso 1 ou 2, mais simples).

\begin{lstlisting}
/* Ex12.22: delete_bst.c -- estrutura central */
TreeNode* deleteNode( TreeNode *root, int valor )
{
    if ( root == NULL ) {
        printf( "Valor %d nao encontrado.\n", valor );
        return NULL;
    }

    if ( valor < root->data ) {
        root->left = deleteNode( root->left, valor );
    } else if ( valor > root->data ) {
        root->right = deleteNode( root->right, valor );
    } else {
        /* Achou o no a remover */

        if ( root->left == NULL && root->right == NULL ) {
            /* Caso 1: folha */
            free( root );
            return NULL;
        }
        else if ( root->left == NULL ) {
            /* Caso 2a: so filho direito */
            TreeNode *temp = root->right;
            free( root );
            return temp;
        }
        else if ( root->right == NULL ) {
            /* Caso 2b: so filho esquerdo */
            TreeNode *temp = root->left;
            free( root );
            return temp;
        }
        else {
            /* Caso 3: dois filhos. Procura substituto:
               maior valor da subarvore esquerda */
            TreeNode *sub = root->left;
            while ( sub->right != NULL ) sub = sub->right;

            root->data = sub->data;
            root->left = deleteNode( root->left, sub->data );
        }
    }
    return root;
}
\end{lstlisting}

\textbf{Observação:} o uso de recursão e retorno de \texttt{TreeNode *} simplifica
drasticamente o código -- evita rastreamento manual do nó pai. Padrão idiomático
em todas as operações em árvores.
\end{respostabox}

\newpage

% ============================================================
\section{Exercício 12.23 -- Busca em Árvore Binária}
% ============================================================

\begin{respostabox}
\begin{lstlisting}
/* Ex12.23: bst_search.c */
TreeNode* binaryTreeSearch( TreeNode *root, int chave )
{
    if ( root == NULL )               return NULL;
    if ( chave == root->data )        return root;
    if ( chave  < root->data )
        return binaryTreeSearch( root->left,  chave );
    return     binaryTreeSearch( root->right, chave );
}
\end{lstlisting}

\textbf{Por que árvores são eficientes para busca?} A cada passo, descartamos
\textit{metade} da árvore restante: $O(\log n)$ em média. Para um milhão de
elementos, isso significa $\sim 20$ comparações, em vez de $\sim 500\,000$ em
uma lista.

Esta é a base de \texttt{std::map} (C++), \texttt{TreeMap} (Java), índices
B-tree em bancos de dados, etc.
\end{respostabox}

% ============================================================
\section{Exercício 12.24 -- Travessia por Nível}
% ============================================================

\begin{respostabox}
\textbf{Algoritmo BFS (Breadth-First Search) com fila:}

\begin{lstlisting}
/* Ex12.24: level_order.c */
void levelOrder( TreeNode *root )
{
    if ( root == NULL ) return;

    /* Fila de ponteiros para no */
    TreeNode *fila[ 1000 ];
    int frente = 0, fim = 0;

    fila[ fim++ ] = root;

    while ( frente < fim ) {
        TreeNode *atual = fila[ frente++ ];
        printf( "%d ", atual->data );

        if ( atual->left  != NULL ) fila[ fim++ ] = atual->left;
        if ( atual->right != NULL ) fila[ fim++ ] = atual->right;
    }
}
\end{lstlisting}

\textbf{Para a árvore do Ex.~12.5, a travessia por nível é:}
\begin{lstlisting}
49  28 83  18 40 71 97  11 19 32 44 69 72 92 99
\end{lstlisting}

\textbf{Aplicações de BFS:}
\begin{itemize}[noitemsep]
  \item Encontrar caminho mais curto em grafos não-ponderados.
  \item Web crawling em largura.
  \item Renderização de árvores de UI.
  \item Algoritmos de inteligência artificial (jogos, planejamento).
\end{itemize}
\end{respostabox}

\newpage

% ============================================================
\section{Exercício 12.25 -- Impressão Visual da Árvore}
% ============================================================

\begin{respostabox}
\textbf{Truque:} é uma travessia in-order \textit{modificada}: visite primeiro
o filho \textbf{direito}, depois imprima o nó atual, depois visite o filho
esquerdo. Como gira a árvore 90° à esquerda, a saída visual é a árvore deitada
com a raiz à esquerda.

\begin{lstlisting}
/* Ex12.25: outputTree.c */
void outputTree( const TreeNode *root, int totalSpaces )
{
    if ( root == NULL ) return;

    /* Visita direita PRIMEIRO (vai para topo da saida) */
    outputTree( root->right, totalSpaces + 5 );

    /* Espacos antes do valor */
    int i;
    for ( i = 1; i <= totalSpaces; ++i ) putchar( ' ' );

    /* Imprime o no */
    printf( "%d\n", root->data );

    /* Visita esquerda DEPOIS (vai para baixo da saida) */
    outputTree( root->left, totalSpaces + 5 );
}
\end{lstlisting}

\textbf{Para a árvore do Ex.~12.5, a saída seria:}
\begin{saida}
\quad\quad\quad\quad\quad\quad\quad\quad\quad\quad\quad\quad\quad\quad 99\\
\quad\quad\quad\quad\quad\quad\quad\quad\quad\quad 97\\
\quad\quad\quad\quad\quad\quad\quad\quad\quad\quad\quad\quad\quad\quad 92\\
\quad\quad\quad\quad\quad 83\\
\quad\quad\quad\quad\quad\quad\quad\quad\quad\quad\quad\quad\quad\quad 72\\
\quad\quad\quad\quad\quad\quad\quad\quad\quad\quad 71\\
\quad\quad\quad\quad\quad\quad\quad\quad\quad\quad\quad\quad\quad\quad 69\\
49\\
\quad\quad\quad\quad\quad\quad\quad\quad\quad\quad\quad\quad\quad\quad 44\\
\quad\quad\quad\quad\quad\quad\quad\quad\quad\quad 40\\
\quad\quad\quad\quad\quad\quad\quad\quad\quad\quad\quad\quad\quad\quad 32\\
\quad\quad\quad\quad\quad 28\\
\quad\quad\quad\quad\quad\quad\quad\quad\quad\quad\quad\quad\quad\quad 19\\
\quad\quad\quad\quad\quad\quad\quad\quad\quad\quad 18\\
\quad\quad\quad\quad\quad\quad\quad\quad\quad\quad\quad\quad\quad\quad 11
\end{saida}

\textbf{Por que isso funciona?} A travessia in-order produz a ordem crescente
(esquerda $\to$ nó $\to$ direita). Trocando a ordem para direita $\to$ nó $\to$
esquerda, produzimos ordem decrescente -- exatamente o que precisamos para
exibir a árvore com os maiores no topo. O recuo (\texttt{totalSpaces}) cresce
com a profundidade, dando o efeito 3D.

\textbf{Aplicações:} esta técnica é a base de visualização de árvores em
ferramentas de \textit{debugging}, editores de XML/JSON, e até no comando Unix
\texttt{tree}.
\end{respostabox}

\subsection*{Exercícios 12.26 a 12.30}

\begin{respostabox}
Os exercícios 12.26 a 12.30 (construção de um compilador para a linguagem Simple,
gerando código SML para o Simpletron) são projetos extensos disponibilizados em
PDF separado pela Deitel \& Associates. Eles consolidam:

\begin{itemize}[noitemsep]
  \item \textbf{12.26}: \textit{lexer} e \textit{parser} para a linguagem Simple
  \item \textbf{12.27}: análise sintática e geração de código SML
  \item \textbf{12.28-30}: otimizações, recursos avançados, testes
\end{itemize}

Sugerimos consultar o documento original no site da Deitel para os enunciados
completos. Estes exercícios são \textit{projetos de final de curso} e exigem
várias semanas de trabalho.
\end{respostabox}

\end{document}
